\subsection{Reduced schemes}
\begin{exer}
	Let $k$ be a field and $P \in k[T_1,\ldots,T_n]$. Then
	$\Spec(k[T_1,\ldots,T_n]/(P))$ is a reduced (resp. irreducible;
	integral) if and only if $P$ is squarefree (resp. admits only one
	irreducible factor; is irreducible).
\end{exer}
\begin{proof}
	dawg
\end{proof}

\begin{exer}
	Let $X$ be a scheme and $x \in X$. Then the image of the morphism
	$\Spec \Sh_{X,x} \to X$ is the set of points of $X$ that specialize to
	$x$.
\end{exer}
\begin{proof}
	It suffices to prove this in the case where $X$ is affine. If $y$ is
	any point of $X$ that specializes to $x$, we have $\Ide_y \subseteq
	\Ide_x$ and so $y$ is in the image of $\Spec \Sh_{X,x}$. The other
	direction is basically the same.
\end{proof}

\begin{exer}
	Let $\Dvr_K$ be a discrete valuation ring with field of fractions $K$
	and uniformizing parameter $t$. Then $\Spec K[T]$ embeds as an open 
	subscheme of $\Spec \Dvr_K[T]$.
\end{exer}
\begin{proof}
	Observe that $(\Dvr_K)_t = K$, so $(\Dvr_K[T])_t = K[T]$. Since
	$\Spec((\Dvr_K[T])_t)$ is homeomorphic to $D(t) \subseteq
	\Spec(\Dvr_K[T])$, this proves the proposition.
\end{proof}

The second part of this exercise is to find the points of $\Spec K[T]$ which
specialize to the point corresponding to the maximal ideal $(T,t)$ of
$\Dvr_K[T]$. To do this, notice that $\Spec K[T]$ corresponds to the ideals $I$
of $\Dvr_K[T]$ such that $I \cap (t) = \{0\}$, i.e. the principal ideals
generated by nonconstant polynomials. Then the only points that specialize to
$(T,t)$ are $(0)$ and $(T)$.

\begin{exer}
	Let $X$ be a scheme having only a finite number of irreducible
	components $\{X_i\}$. Then $X$ is connected if and only if for any pair
	$i,j$, there exist indices $i_0 = i, i_1,\ldots,i_r = j$ such that
	$X_{i_l} = X_{i_{l+1}}$ for every $l < r$. Furthermore, $X$ is integral
	if and only if it is connected and integral at every point.
\end{exer}
\begin{proof}
	Since the irreducible components of $X$ are finite in number, they must
	all be open; it is clear that having $X$ connected implies the
	condition. On the other hand, suppose that $X$ has the property. If $U$
	is a nonvoid open set such that $\overline{U} = U$, we have that $U$ i
	the union of irreducible components of $X$, so suppose that $X_i$ is any
	component (not neccessarily contained in $U$). Then we have a sequence
	$X = X_0, \ldots, X_r \subseteq U$ such that $X_i \cap X_{i+1} \neq
	\emptyset$. However, since $X_{r-1} \cap U \neq \emptyset$, $X_{r-1}$ is
	contained in $U$ and so we may drop the last index. Repeating this, we
	may assume that $r = 0$ and so $X \subseteq U$.\\

	Now, if $X$ is integral, it is obvious that it is connected and integral
	at every point, so assume the latter. At any point $x \in X$ let 
	$i: \Spec \Sh_{X,x} \to X$ be the canonical morphism. Since
	$i^{-1}(i((0))) \subseteq i^{-1}(\overline{i((0))})$, we have that
	$\Spec \Sh_{X,x} \subseteq i^{-1}(\overline{i((0))})$ and so the image
	of $i$ has a closed point. It follows from exercise 2 that each point
	$x$ is contained in exactly one irreducible component, so the
	connectedness hypothesis implies that $X$ is irreducible. That $X$ is
	reduced is obvious, so $X$ is integral.
\end{proof}

\begin{exer}
	Let $X$ be a scheme. Every irreducible component of $X$ is contained in
	a connected component. If $X$ is locally noetherian, then the connected
	components are open. If $X$ is noetherian, the connected components are
	finite in number.
\end{exer}
\begin{proof}
	The first and third statements are clear, and the second one follows
	from the third.
\end{proof}

\begin{exer}
	Let $A$ be a ring and $X$ be a scheme.
	\begin{enumerate}[label=(\alph*)]
		\item The following properties are equivalent:
			\begin{enumerate}[label=(\roman*)]
				\item $X$ is connected;
				\item $\Sh_X(X)$ has no idempotent elements
					other than $1$;
				\item $\Spec(\Sh_X(X))$ is connected.
			\end{enumerate}
		\item Any local scheme is connected.
		\item Suppose that the connected components of $X$ are open. Let
			$U$ be a connected component of $X$. Then there is a
			unique idempotent element $e$ of $\Sh_X(X)$ such that
			${e|}_U = 1$ and ${e|}_{X \setminus U} = 0$. This
			establishes a canonical bijection from the set of
			connected components of $X$ onto the set of
			indecomposable idempotent elements of $\Sh_X(X)$, the
			inverse map being given by $e \mapsto V(1-e)$.
	\end{enumerate}
\end{exer}
\begin{proof}
	For part (a), let $e \neq 1 \in \Sh_X(X)$ be an idempotent element. At
	any point $x \in X$, $e_x$ must also be idempotent, so generates an
	ideal $(e_x)$ that squares to itself. If $e_x \neq 1$, Nakayama's lemma
	implies that it is $0$. We see that the sets $U = \{x \in X: e_x = 1\}$
	and $V = \{x \in X: e_x = 0\}$ are disjoint and form an open cover of
	$X$, so (i) implies (ii).\\

	On the other hand, suppose that $U,V$ are disjoint open sets that cover
	$\Spec(\Sh_X(X))$. Then there is an element of $\Sh_X(X)$ that restricts
	to $0$ on $U$ and $1$ on $V$. It is clearly idempotent and not equal to
	$1$, so (ii) implies (iii).\\

	Finally, suppose that $X$ is disconnected. Then the same argument as in
	part (ii) implies that there is an idempotent $e \neq 1 \in \Sh_X(X)$,
	hence $\Spec(\Sh_X(X))$ is disconnected and so (a) follows.\\

	Part (b) follows from the same reasoning as part (a)(ii), and part (c)
	is obvious.
\end{proof}
